%% ch-meta-soundness.tex -- Part 2 (\ref{part:meta}), the soundness of the core for version-X (link 3).
%% Hand-written (DS-6).  Cites the design specification as \DS{n}.  Placed after ch-meta-k2.
\providecommand{\DS}[1]{\hyperref[ds:#1]{DS-#1}}
\providecommand{\Sq}{\Rightarrow}
\providecommand{\FV}{\mathrm{FV}}

\chapter{Soundness of the core}
\label{ch:meta-soundness}

This chapter is link~3 of Section~\ref{sec:meta-chain}.  It argues
that each of the 56 core rules of Section~\ref{sec:meta-core} is sound
for version-X, and that each axiom of the catalog holds, under one
reading of VNB's terms and formulas.

The argument is written by hand (\DS{6}).  It is not a VNB theorem
(Section~\ref{sec:meta-not-claimed}).  The functoid part is
Chapter~\ref{ch:meta-functoids}; here every formula is read after the
elimination of its defined heads.

%% =========================================================
\section{The reading}\label{sec:meta-reading}

\subsection{Models}

Let $M$ be a model of version-X, NBG with global choice
(\DS{3})~\cite{vonNeumann1925,Bernays1937,Godel1940}.  The objects of
$M$ are its classes; a set is a class that is a member of a class.  Fix
in $M$ a well-ordering $W$ of the universe.  NBG with global choice
provides one.

A VNB class is read as a class of $M$ (\DS{4}).  Free names
(variables) range over all classes of $M$.  An \emph{assignment}
$\sigma$ gives each free name a class.

\subsection{Terms}

A term denotes a class of $M$ or nothing.  We write $[\![t]\!]\sigma$
for the denotation, and say that $t$ \emph{denotes} when there is one.
The clauses:
\begin{itemize}
\item A free name denotes $\sigma$ of it.  A primitive constant denotes
  a fixed class: $\mathrm{SET}$ the class of all sets, $\mathrm{ORD}$
  the class of von Neumann ordinals, $\mathrm{NN}$ the von Neumann
  naturals, $\mathrm{EMPTY\text{-}SET}$ the empty set, the numerals and
  the number classes as in Section~\ref{sec:meta-numbers}.
\item \emph{Lists.}  The list $[e_1..e_n]$ denotes when every $e_i$
  denotes a set (\DS{24}).  Its value is the Kuratowski pair
  $(n, s)$, where $s$ is the set of Kuratowski pairs $(i, e_i)$,
  $1\leq i\leq n$.  These pairs are NBG's own, not VNB's; a VNB pair
  $[x, y]$ is the list of length 2.  Every list is a set.
\item $\mathrm{LENGTH}(c)$ is $n$ when $c$ is a list of length $n$, and
  $0$ for every other class.  $\mathrm{NTH}(k, c)$ is the $k$-th entry
  when $c$ is a list and $1\leq k\leq n$, and nothing otherwise.
  $\mathrm{CONS}(x, c)$, for a set $x$, is the list with first entry $x$
  followed by the entries of $c$ when $c$ is a list, and $[x]$
  otherwise; for a proper class $x$ it is nothing.
\item The class constructors $\mathrm{UNION}$, $\mathrm{INTERSECTION}$,
  $\mathrm{CARTESIAN}$, $\mathrm{TUPLES}$ and $\mathrm{FUN}$ denote
  when every argument denotes.  They are then the classes of sets that
  their membership schemas describe.  $\mathrm{CARTESIAN}(A_1..A_n)$ is the class of lists of
  length $n$ whose $i$-th entry is in $A_i$; $\mathrm{TUPLES}(A)$ the
  class of lists all of whose entries are in $A$; $\mathrm{FUN}(A)$ the
  class of sets $f$ of VNB pairs, single-valued, with domain exactly
  $A$; $\mathrm{FUN}(A, B)$ those members of $\mathrm{FUN}(A)$ whose
  values lie in $B$.  $\mathrm{UNION}()$ is the empty set; the nullary
  intersection and product are fixed arbitrarily, since no rule or
  axiom mentions them.
\item $\mathrm{succ}$ and $\mathrm{succ\_ORD}$ are total: at a set $x$
  the value is $x\cup\{x\}$, at a proper class the class itself.
\item \emph{Application.}  $\mathrm{APPLY}(f, x)$ is the unique $v$ with
  $[x, v]\in f$, when $f$ and $x$ denote and exactly one such $v$
  exists, and nothing otherwise.  Such a $v$ is a set.  An application
  of arity $n\geq 2$, $f(a_1,\ldots,a_n)$, is
  $\mathrm{APPLY}(f, [a_1..a_n])$.
\item $\mathrm{IF}(p, a, b)$ is the denotation of $a$ when $p$ is true,
  of $b$ otherwise, and nothing when that branch denotes nothing.
\item $\iota x.\,B$ denotes the unique class $C$ such that $B$ holds at
  $x := C$, and nothing when there is no such class or more than one.
\item $\mathrm{CHOICE}(A)$ denotes the $W$-least member of $A$ when $A$
  denotes an inhabited class, and nothing otherwise (\DS{52}).
\item \emph{Binder terms.}  Over a denoting domain $A$:
  $\mathrm{SEP}\,x{\in}A.\,p$ is the class of the members of $A$ at which
  $p$ holds; $\bigcup x{\in}A.\,b$ the class of the members of the
  denoting values $b[x]$, $x\in A$; $\lambda x{\in}A.\,b$ the class of
  the pairs $[x, b[x]]$, $x\in A$, $b[x]$ a set.
  $\mathrm{COMP}\,x.\,p$ is the class of the sets at which $p$ holds.
  Each denotes when that class exists in $M$, and nothing otherwise.
  Over a non-denoting domain the binder term denotes nothing.
\end{itemize}
Each other constructor is read by its defining axioms in the catalog
(Section~\ref{sec:meta-axioms}).  A guarded definition is read as a
description: off every guard its head denotes nothing (\DS{31}).

\subsection{Formulas and sequents}

The atoms $\in$, $=$, $\leq$ and $<$ are strict: an atom with a
non-denoting argument is false (\DS{18}).  So $t = t$ says that $t$
denotes.  $s\equiv t$ holds when neither side denotes, or both denote the
same class.  A predicate application $P(t_1..t_n)$ is false when some
$t_i$ does not denote, and otherwise holds when the values satisfy $P$.
Primitive predicates are classes of $M$; a defined predicate is read
through its definition.

The connectives are classical.  The quantifiers range over all classes
of $M$.  A formula containing defined heads means its elimination
(\DS{7}; Chapter~\ref{ch:meta-functoids}).

A sequent $\Gamma\Sq A$ is \emph{valid in $M$} when every assignment
that makes each member of $\Gamma$ true makes $A$ true.  A theory $T$ is
\emph{true in $M$} when each of its members is true at every
assignment.

\begin{dfn}\label{dfn:meta-sound}
A core rule $\mathrm{CR}\text{-}x$ is \emph{sound} when, for every model
$M$ of version-X and every theory $T$ true in $M$, each instance of
$\mathrm{CR}\text{-}x(T, \mathit{hs}, c)$ whose hypotheses are valid in
$M$ has a conclusion valid in $M$.
\end{dfn}

A definedness side sequent $\Gamma\Sq t=t$ is an ordinary hypothesis
under this definition.  Nothing else is needed for it.

\subsection{Two lemmas}

\paragraph{Substitution.}  Let $t$ be a closed term that denotes at
$\sigma$.  Then $B[x:=t]$ is true at $\sigma$ exactly when $B$ is true
at $\sigma$ with $x$ given the value of $t$.  The same holds for the
denotation of a term.  If $s\equiv t$ holds at $\sigma$, then $P[z:=s]$
and $P[z:=t]$ have the same truth value at $\sigma$.  Both statements
are proven by induction on the expression.  They hold because each clause
above reads a subterm only through its denotation or the absence of one.
The encoding is locally nameless (\DS{10}), so no capture arises.

\paragraph{Elimination.}  Let $p$ satisfy $\mathrm{NBG\text{-}OK}$: every
quantifier and description in it is confined to sets, read through the
tables $\mathrm{FUN\text{-}DEF}$ and $\mathrm{PRED\text{-}DEF}$ (\DS{39}).
Then $p$ is equivalent in $M$ to a formula of NBG all of whose
quantifiers range over sets, with class parameters: the values of its
free names, the primitive classes and $W$.

The proof replaces each atom about a compound term by its defining
condition.  $y\in\mathrm{SEP}\,x{\in}A.\,q$ becomes
$y\in A\wedge q[x:=y]$.  $y=\mathrm{APPLY}(f,x)$ becomes
$[x,y]\in f\wedge\forall y'\,([x,y']\in f\to y'=y)$, where $y'$ is
confined by the membership.  $y=\mathrm{CHOICE}(A)$ becomes
$y\in A\wedge\forall z\,(z\in A\to\neg\,zWy)$.  A list becomes its
Kuratowski code.  A confined description becomes the confined unique
existence it abbreviates.  Hence the class $\{x\in\mathrm{SET} : p\}$
exists in $M$, by the class comprehension of NBG.

\begin{prop}[hand-written, conditional]\label{prop:meta-sound}
Assume the hypotheses of Section~\ref{sec:meta-sound-hyp}.  Every core
rule is sound.  Consequently, if $T$ is true in a model of version-X,
every sequent in $\mathrm{DERIVABLE}(T)$ is valid in it, and
$\{\}\Sq\textsc{falsity}$ is not derivable.
\end{prop}

The consequence is an induction on $\mathrm{DERIVABLE}(T)$, the least
set closed under the rules.  $\{\}\Sq\textsc{falsity}$ is valid in no
model.  If version-X is consistent it has a model, so $T$, when true in
one, derives no \textsc{falsity}.  Section~\ref{sec:meta-axioms} argues
that the catalog is true in every model of version-X, with the primitive
heads read as stated there.  A theorem enters $T$ only after it is derived
(links~1 and~2), so by induction on the order of installation it is true
as well.

%% =========================================================
\section{The rules}\label{sec:meta-rules}

The groups follow the table of Section~\ref{sec:meta-core}.  In each
argument, $\sigma$ is an assignment that makes $\Gamma$ true, and a
hypothesis is used through its validity at $\sigma$.

\subsection{Logic (21 rules)}

and-intro, and-elim, or-intro-left, or-intro-right, or-elim,
implies-intro, backchain, not-intro, not-elim, iff-intro, iff-elim,
assumption, cut, weakening and proof-by-contradiction are the
propositional rules of classical natural deduction in the sequent
form~\cite{Gentzen1935}.  Each is sound by the truth tables.  The
elimination rules read a formula of the context; that changes nothing in
the argument.  truth-intro concludes \textsc{truth},
true everywhere.  theorem-assumption concludes a member of $T$, true by
the assumption on $T$.

The four quantifier rules are those of a partial first-order logic in
which variables always denote and a term may not, as in the logic of
IMPS~\cite{IMPS} and in~\cite{Farmer2023}.
\begin{itemize}
\item forall-elim and forsome-intro carry the hypothesis $\Gamma\Sq t=t$
  (\DS{19}, \DS{40}).  At $\sigma$, $t$ denotes a class $C$.  By the
  substitution lemma, $B[x:=t]$ is $B$ at $x:=C$.  The universal gives
  it; the existential is witnessed by it.  Without the hypothesis
  neither rule is sound: $\forall x.\,x=x$ is true, and
  $\iota x.\,\textsc{falsity} = \iota x.\,\textsc{falsity}$ is false.
\item forall-intro: $y\notin\FV(\Gamma)\cup\FV(B)$.  For a class $C$,
  the assignment $\sigma[y:=C]$ still makes $\Gamma$ true, so $B[x:=y]$
  holds there, which is $B$ at $x:=C$.  forsome-elim is the dual
  argument, with $y\notin\FV(G)$.
\end{itemize}

\subsection{Equality (4 rules)}

reflexivity concludes $t=t$ for $t$ a free name or an atomic constant.
A free name denotes under every assignment.  An atomic constant denotes
by hypothesis~(S2).  quasi-reflexivity holds by the definition of
$\equiv$.

eq-subst and eq-subst-sym need one of $s=t$, $s\equiv t$, $t=s$,
$t\equiv s$ in $\Gamma$ (\DS{41}).  Each makes $s\equiv t$ true at
$\sigma$.  The substitution lemma then gives $P[z:=s]$ from $P[z:=t]$.
The rule is sound for $\equiv$, not only for $=$, because the reading
is compositional in the denotation or its absence.

\subsection{Structural membership and tuples (7 rules)}

Each schema is the membership condition of the reading.  The work is at
the arguments that may not denote.
\begin{itemize}
\item union-mem, with one hypothesis $\Gamma\Sq S_i=S_i$ per entry
  (\DS{53}).  All entries denote, so $\bigcup(S_1..S_n)$ is the union of
  their values, and the equivalence is its definition.  If $x$ does not
  denote, both sides are false.  At $n=0$ the union is empty and the
  right side is \textsc{falsity}.  Without the hypotheses the schema is
  false: with $S_2$ non-denoting and $x\in S_1$ true, the right side holds
  and the left does not, since $\mathrm{UNION}$ is strict.
\item intersection-mem, $n\geq 1$.  If some $S_i$ does not denote, both
  sides are false.  Otherwise the equivalence is the definition.  No
  hypothesis is needed, since every entry occurs in a strict atom on the
  right.
\item cartesian-mem, $n\geq1$.  The same argument.  The right side's
  conjuncts $\mathrm{NTH}(i,x)\in A_i$ are false at a non-denoting
  $A_i$, and the product does not denote.
\item tuples-mem, with the hypothesis $\Gamma\Sq A=A$ (\DS{53}).  If some
  $e_i$ is a proper class or does not denote, the list does not denote
  and $e_i\in A$ is false: both sides are false.  Otherwise the list is
  in $\mathrm{TUPLES}(A)$ exactly when each entry is in $A$.  At $n=0$
  the list $[\,]$ denotes and lies in $\mathrm{TUPLES}(A)$ for every
  denoting $A$; the right side is \textsc{truth}.  This is where the
  hypothesis is needed.
\item nth-eq and length-eq carry $\Gamma\Sq L=L$ for the literal
  $L=[e_1..e_n]$ (\DS{25}).  $L$ denotes, so each $e_k$ is a set,
  $\mathrm{NTH}(k,L)$ is $e_k$ for $1\leq k\leq n$, and
  $\mathrm{LENGTH}(L)$ is $n$.  The conclusion is the quasi-equation.
\item tuple-eq carries $\Gamma\Sq L=L$ and $\Gamma\Sq M=M$.  Two
  denoting lists of the same length are equal exactly when their entries
  are, by the uniqueness of Kuratowski pairs.  This is also the claim of
  the primitive \texttt{tuple-equality-decompose} (\DS{50}).
\end{itemize}

\subsection{Conditional terms (2 rules)}

The rule if-true has the hypotheses $\Gamma\Sq p$ and
$\Gamma,\ \mathrm{IF}(p,a,b)\equiv a\Sq G$, and the conclusion
$\Gamma\Sq G$.  At $\sigma$, $p$ is true, so $\mathrm{IF}(p,a,b)$
has the denotation of $a$, or none when $a$ has none.  The added
quasi-equation is therefore true at $\sigma$, and the second hypothesis
gives $G$.  if-false is the same with $\neg p$ and $b$.  The core
concludes the quasi-equation, not the strict one: the strict equation
would assert that the branch denotes (\DS{20}).

\subsection{Set formation, descriptions, functions (9 rules)}

Every schema here that forms a class carries $\mathrm{NBG\text{-}OK}$ on
its condition or body.  By the elimination lemma the class it names
exists in $M$.  This is the only place where the strength of NBG, and
not of Morse--Kelley, is at stake (\DS{38}, \DS{39}).
\begin{itemize}
\item sep-mem and comp-mem are the definitions of the two classes.  At a
  non-denoting $t$ both sides are false: $t\in A$, $t\in\mathrm{SET}$
  are strict.
\item sep-sethood: a subclass of a set is a set (separation in NBG).
\item big-union-mem is the definition of the union of the denoting values.
  $\exists z\,(z\in A\wedge x\in b)$ is confined by $z\in A$.
\item big-union-sethood: $A$ is a set and each value $b[z]$, $z\in A$, is
  a set.  By replacement the values form a set; its union is a set.
\item iota-def: the first hypothesis says that exactly one class
  satisfies $B$.  So $\iota x.\,B$ denotes it, and $B[x:=\iota x.\,B]$ is
  true by the substitution lemma.  The second hypothesis gives $G$.
\item iota-in-elim: $\iota x.\,B$ denotes, so it denotes the unique
  class satisfying $B$, and $B[x:=\iota x.\,B]$ holds.
\item lambda-type carries $\Gamma\Sq\forall x.\,(x\in A\to b\in D)$,
  $\Gamma\Sq A\in\mathrm{SET}$ and $\Gamma\Sq D=D$ (\DS{53}).  Let $L$ be
  $\lambda x{\in}A.\,b$.  $L$ is a class, by the elimination lemma.  Each
  value $b[x]$, $x\in A$, is in $D$, so it is a set, and $L$ contains a
  pair for every $x\in A$.  $L$ is single-valued with domain $A$.  It is a
  set by replacement, since $A$ is a set.  $D$ denotes, so
  $\mathrm{FUN}(A,D)$ denotes, and $L$ is a member.  Without $D=D$ the
  instance $A=\mathrm{EMPTY\text{-}SET}$, $D$ non-denoting has a true
  first hypothesis and a false conclusion.
\item lambda-beta-eq carries $\Gamma\Sq t\in A$ and
  $\Gamma\Sq b[x:=t]\in\mathrm{SET}$.  $t$ and $b[t]$ are sets, so
  $[t, b[t]]$ denotes and lies in $L$.  A pair $[x, v]\in L$ with $x=t$
  has $v=b[t]$.  So $\mathrm{APPLY}(L,t)$ is $b[t]$, and the
  quasi-equation holds.  Without the second hypothesis the rule is
  unsound: a pair whose second entry is a proper class does not denote,
  so a lambda term with a class-valued body has no pair at that point,
  and the application denotes nothing while $b[t]$ denotes.
\end{itemize}

\subsection{Definedness (9 rules; \DS{42})}

\begin{itemize}
\item def-strict-in-l, def-strict-in-r: a true membership has denoting
  sides.
\item def-total: $\mathrm{UNION}$, $\mathrm{INTERSECTION}$,
  $\mathrm{CARTESIAN}$, $\mathrm{TUPLES}$, $\mathrm{FUN}$, $\mathrm{succ}$
  and $\mathrm{succ\_ORD}$ denote at denoting arguments.  For the class
  constructors this is comprehension over a predicative condition of the
  arguments.  For the two successors it is their totality.
\item def-binder: $\mathrm{SEP}$, $\bigcup$ and $\lambda$ over a
  denoting domain, with an $\mathrm{NBG\text{-}OK}$ body, denote by the
  elimination lemma.  def-comp: likewise for $\mathrm{COMP}\,x.\,p$.
\item def-list: a list of sets denotes (\DS{24}).
\item def-strict-arg: the heads $\mathrm{UNION}$ through $\mathrm{FUN}$
  (with $\mathrm{NTH}$, $\mathrm{LENGTH}$, $\mathrm{LIST}$ and
  $\mathrm{APPLY}$), $\mathrm{succ}$ and $\mathrm{succ\_ORD}$ denote only
  at denoting arguments.  For $\mathrm{LIST}$ this is the other half of
  \DS{24}.  def-strict-dom: a binder term denotes only over a denoting
  domain.
\item def-strict-pred: a true predicate application has denoting
  arguments, by the reading.  The reading agrees with the definition of
  a defined predicate $P$ only at denoting arguments.  This is enough,
  because every instantiation of a definition owes the definedness of the
  instance (\DS{19}): a definition is never applied at a term that may
  not denote.
\end{itemize}

\subsection{Congruence, functoids, tuples (4 rules)}

\begin{itemize}
\item iota-cr: from $\Gamma\Sq B[x:=y]\leftrightarrow C[x:=y]$, $y$ fresh,
  infer $\iota x.\,B\equiv\iota x.\,C$.  By forall-intro's argument $B$ and
  $C$ hold at the same classes.  So the two descriptions denote the same
  class, or both denote nothing.
\item tbind-cr: for $\mathrm{SEP}$ and $\mathrm{COMP}$ the premise is
  the equivalence of the conditions, for $\bigcup$ and $\lambda$ the
  quasi-equation of the bodies, at a fresh $y$.  The two classes have the
  same members, or neither exists (\DS{43}).  No predicativity premise is
  needed: the conclusion is a quasi-equation, true when neither side
  denotes.
\item functoid-beta-eq: $F(l)\equiv b[l]$ for a row $(F,|l|,b)$ of
  $\mathrm{FUN\text{-}DEF}$.  The meaning of $F(l)$ is that of its
  instantiated body (\DS{7}, \DS{34}), at every argument list, so the
  quasi-equation holds.  It rests on hypothesis~(S3).
\item tuple-eta, $n\geq1$: a member $x$ of $\mathrm{CART}(A_1..A_n)$ is a
  list of length $n$ of sets, and a list is the list of its own entries.
  The right side denotes, being equal to $x$ (\DS{45}).
\end{itemize}

%% =========================================================
\section{The axioms}\label{sec:meta-axioms}

The catalog is the set of axioms in \texttt{reference/THEOREMS.md},
section ``Axioms -- the fixed VNB base''.  Its sources are the base
theory of \texttt{library.scm} and the primitive shelf:
\texttt{theorem-library/axioms.scm}, \texttt{number-systems.scm},
and, in \texttt{structure-library/}, \texttt{ordinals.scm},
\texttt{list-recursion.scm}, \texttt{injection.scm},
\texttt{extended-reals.scm}, \texttt{order-lemmas.scm} and
\texttt{ca-lft-vocab.scm}.  Each
axiom falls into one of the classes below.

\subsection{NBG axioms and their direct consequences}

\begin{itemize}
\item Extensionality: \texttt{class-extensionality},
  \texttt{extensionality}.
\item Sets: \texttt{membership-implies-sethood},
  \texttt{empty-set-is-set}, \texttt{empty-set-has-no-members},
  \texttt{subset-set}.
\item Pairing, power and union: \texttt{pairing},
  \texttt{pairing-membership}, \texttt{power-set},
  \texttt{power-set-membership}, \texttt{union-set-closure}.
\item Separation: \texttt{intersection-set-closure} (one factor a set),
  \texttt{complement-in-set-closure}.
\item Replacement: \texttt{image-set}.  Here $\mathrm{IMAGE}(\varphi, S)$
  is the class of sets $w$ with $\varphi(x)=w$ for some $x\in S$
  (\texttt{image-membership-iff}).  This is the range of the
  single-valued part of $\varphi$ over $S$, a set when $S$ is one.
  $\varphi$ ranges over classes: a bare functoid is never an instance
  (\DS{33}).
\item Infinity: \texttt{nn-is-set}.
\item Global choice: \texttt{choice-axiom}, by the reading of
  $\mathrm{CHOICE}$ (\DS{52}).
\item Theorems of NBG: \texttt{burali-forti}; \texttt{nn-induction} and
  \texttt{transfinite-induction}, at every class $c$.  Every non-empty
  class of ordinals has a least member.
\item Equality.  The axioms \texttt{equality-symmetry},
  \texttt{equality-transitivity} and \texttt{eq-subst-membership} are
  true for strict equality.
\end{itemize}

\subsection{Definitional axioms}

Each of these fixes the meaning of a head under the reading.  It is
therefore true by choice of that meaning.  A membership condition is
predicative, read through defined heads, as the binder-predicativity
gate checks for every definitional or primitive membership equivalence
(\DS{38}).  So the class it describes exists.
\begin{itemize}
\item Application and domain: \texttt{app-graph} (the reading of
  $\mathrm{APPLY}$), \texttt{dom-membership}, \texttt{is-fun-def}, and
  the definitional \texttt{apply-tupling-2} and \texttt{-3}, the
  application of arity $n\geq2$.
\item Functions: \texttt{fun-no-junk}, \texttt{fun-codomain-iff},
  \texttt{fun-domain-apply-def}.  These are the reading of
  $\mathrm{FUN}$.  At $x\notin A$ no pair $[x, v]$ is in $f$, so $f(x)$
  denotes nothing.
\item Boolean operations.  The axioms \texttt{union-membership},
  \texttt{intersection-membership},
  \texttt{complement-in-membership} and \texttt{subset-def}; and the
  schemas \texttt{union-decompose},
  \texttt{intersection-decompose}.  The two schemas are true at every
  instance whose entries denote.  The macete door owes each entry's
  definedness (\DS{19}), and the nullary intersection is refused.
\item Constructions: \texttt{make-set-membership},
  \texttt{make-set-empty}, \texttt{partial-fun-membership},
  \texttt{partial-fun-binary-membership}, \texttt{res-apply},
  \texttt{res-typing}, \texttt{res-codomain}.  $\mathrm{RES}(f,b)$ is the
  set of the pairs of $f$ whose first entry is in $b$.
\item Quasi-equality: \texttt{quasi-eq-reflexivity}, \texttt{-symmetry},
  \texttt{-transitivity}, \texttt{quasi-eq-def},
  \texttt{quasi-eq-subst-membership}.  The last holds because $b\in S$
  makes $b$ denote, and then $a\equiv b$ is $a=b$.
\item Ordinals: $\mathrm{ORD\text{-}LE}$ is inclusion on $\mathrm{ORD}$
  and false elsewhere; $\mathrm{succ\_ORD}(\alpha)=\alpha\cup\{\alpha\}$;
  $\mathrm{ORD\text{-}SEGMENT}(\alpha)=\alpha$;
  $\mathrm{SUP\text{-}ORD}(a)=\bigcup a$.  Under these, \texttt{ord-le-*},
  \texttt{ord-succ-*}, \texttt{ord-segment-*}, \texttt{sup-ord-*},
  \texttt{ord-zero-least}, \texttt{limit-ord-is-sup},
  \texttt{nn-subset-ord}, \texttt{ord-succ-nn} and
  \texttt{ord-le-nn-compat} are facts about von Neumann ordinals.
  \texttt{ord-lt-iff} and \texttt{limit-ord-iff} define two predicates.
\item Structures.  For each declared structure $X$, $\mathrm{IS\text{-}}X$
  is a predicate defined by its shape and laws.  Its last conjunct is
  $s\in\mathrm{SET}$, and a refinement inherits it through its parent.
  The class axiom $\forall s\,(s\in X\leftrightarrow\mathrm{IS\text{-}}X(s))$
  is then comprehension over that condition.  It holds because the
  condition implies sethood and is predicative (the gate above).  The
  morphism and hom-set definitions are definitions of new heads.
\end{itemize}

\subsection{The two sethood axioms}

\texttt{fun-set-iff} and \texttt{cartesian-set-iff} are implications.
Their names are kept:
\begin{itemize}
\item \texttt{fun-set-iff}:
  \[\forall A\,\forall B\,(A\in\mathrm{SET}\wedge B\in\mathrm{SET}\to\mathrm{FUN}(A,B)\in\mathrm{SET});\]
\item \texttt{cartesian-set-iff}:
  \[\forall A\,\forall B\,(A\in\mathrm{SET}\wedge B\in\mathrm{SET}\to\mathrm{CARTESIAN}(A,B)\in\mathrm{SET}).\]
\end{itemize}
\emph{cartesian-set-iff.}  The map $(a,b)\mapsto[a,b]$ is a class
function on the NBG product $A\times B$, a set.  By replacement its
range is a set, and that range is $\mathrm{CARTESIAN}(A,B)$.

\emph{fun-set-iff.}  A member of $\mathrm{FUN}(A,B)$ is a set of pairs
$[x,y]$ with $x\in A$ and $y\in B$, so a subset of
$\mathrm{CARTESIAN}(A,B)$.  Hence $\mathrm{FUN}(A,B)$ is a subclass of
$\mathrm{POWER}(\mathrm{CARTESIAN}(A,B))$, a set.

The converses do not hold in general.  $\mathrm{CARTESIAN}(\mathrm{EMPTY\text{-}SET},\mathrm{ORD})$
is empty, hence a set.  $\mathrm{FUN}(\mathrm{EMPTY\text{-}SET},\mathrm{ORD})$
is $\{\mathrm{EMPTY\text{-}SET}\}$.  $\mathrm{FUN}(\mathrm{ORD},B)$ is
empty for every $B$, since a set function has a set domain.  The
converses that hold are theorems, proven in
\texttt{theorem-library/set-iff-converses.scm}:
\begin{itemize}
\item \texttt{fun-domain-in-set}: an inhabited $\mathrm{FUN}(A,B)$ has a
  set domain $A$;
\item \texttt{fun-set-converse}: if $A$ is an inhabited set and
  $\mathrm{FUN}(A,B)$ is a set, then $B$ is a set;
\item \texttt{cartesian-set-converse-1} and \texttt{-2}: if
  $\mathrm{CARTESIAN}(A,B)$ is a set and one factor is inhabited, the
  other factor is a set.
\end{itemize}

\subsection{Consequences under the reading}

These follow from the reading by a short argument.
\begin{itemize}
\item Functions.  The axioms \texttt{dom-of-fun},
  \texttt{dom-fun-membership}, \texttt{fun-apply-type},
  \texttt{fun-elements-are-sets},
  \texttt{fun-domain-elements-are-sets},
  \texttt{fun-domain-extensionality} and
  \texttt{partial-fun-binary-sethood}.  The last is the union, over the subsets $b$ of $A$, of the sets
  $\mathrm{FUN}(b,C)$.
\item Lists: \texttt{list-sethood}, \texttt{tuples-sethood},
  \texttt{length-in-nn}, \texttt{length-of-empty}, \texttt{nth-in-range},
  \texttt{tuple-members-are-sets}, \texttt{make-set-sethood}.  A list's
  entries are sets.  The lists over a set $A$ form a set: they are coded
  by finite sequences from $A$.
\item List generation, primitive: \texttt{tuple-length-zero} and
  \texttt{tuple-cons-decompose}.  Under the reading every list of length
  $n+1$ is $\mathrm{CONS}$ of its first entry onto the list of the
  others.
\item List construction, definitional: \texttt{empty-in-tuples},
  \texttt{cons-in-tuples}, \texttt{length-cons}, \texttt{nth-cons-1},
  \texttt{nth-cons-succ}, \texttt{makeset-cons}.  The equations are
  unguarded in $L$.  They hold at a non-list $L$ because
  $\mathrm{LENGTH}(L)$ is $0$ and $\mathrm{CONS}(x,L)$ is $[x]$.  At that
  $L$, $\mathrm{NTH}$ has no index in range and $\mathrm{MAKE\text{-}SET}(L)$
  is empty.  A reading in which $\mathrm{LENGTH}$ or $\mathrm{CONS}$
  denotes nothing off the lists would make \texttt{length-cons} and
  \texttt{nth-cons-1} false.
\item \texttt{nn-not-a-list} (\DS{49}).  Every list is a Kuratowski pair
  $\{\{n\},\{n,s\}\}$, and no member of a Kuratowski pair is empty.
  $0$ is empty, and every natural $n\geq1$ has $0$ as a member.
  So no natural is a list.
\item The primitive rule \texttt{cartesian-decompose} is tuple eta with
  the membership condition, and \texttt{tuple-equality-decompose} is
  tuple-eq (\DS{50}).
\end{itemize}

\subsection{The number systems}\label{sec:meta-numbers}

The axioms of \texttt{number-systems.scm}, the order and infinity
primitives and \texttt{recip-defined-nonzero} hold in a standard
construction, recoded so that the inclusions are literal.

Build $\mathbb{C}$ over the Dedekind reals in $M$.  Recode it by an
injection into the sets that is the identity from the standard copy of
$\omega$ onto the von Neumann naturals.  Then
$\mathrm{NN}\subseteq\mathrm{ZZ}\subseteq\mathrm{QQ}\subseteq\mathrm{RR}\subseteq\mathrm{CC}$
holds literally, $0$ is the empty set, and $\mathrm{succ}$ on
$\mathrm{NN}$ is $n+1$.  The arithmetic heads, the order, $|\cdot|$,
$\max$, $\min$, conjugation and the powers are transported.  Then:
\begin{itemize}
\item the field, order and completeness axioms (\texttt{nn-*},
  \texttt{zz-*}, \texttt{qq-*}, \texttt{rr-*}, \texttt{cc-*},
  \texttt{power-*}) hold, \texttt{rr-sup-*} included;
\item $\mathrm{recip}$ is the reciprocal on the non-zero numbers and
  denotes nothing at $0$, which gives \texttt{recip-defined-nonzero};
\item $\mathrm{RR\text{-}STAR}$ is $\mathrm{RR}$ with two fresh sets
  $\mathrm{NEG\text{-}INF}$ and $\mathrm{POS\text{-}INF}$.  $\leq$ is the
  order of the extended line on $\mathrm{RR\text{-}STAR}$ and false
  elsewhere.  This gives \texttt{pos-inf-upper-bound},
  \texttt{neg-inf-lower-bound}, \texttt{rr-star-leq-antisymmetric} and
  \texttt{rr-star-leq-domain}.  The two points must be sets for the
  definitional \texttt{rr-star-membership} to hold;
\item $\mathrm{CC\text{-}INF}$ is a fresh set outside $\mathrm{CC}$
  (\texttt{cc-inf-in-set}, \texttt{cc-inf-not-in-cc});
\item the definitional \texttt{binary-minus-def},
  \texttt{binary-divide-def}, \texttt{rr-abs-def}, \texttt{rr-max-def}
  and \texttt{rr-min-def} define new heads.
\end{itemize}
The recoding does not disturb the reading of lists: no axiom says
whether a number other than a natural is a list.

%% =========================================================
\section{The hypotheses}\label{sec:meta-sound-hyp}

Proposition~\ref{prop:meta-sound} assumes the following.
\begin{enumerate}
\item[(S1)] Version-X is consistent.  The oracles are correct
  (\DS{5}).
\item[(S2)] Every atomic constant code of the encoding names a
  denoting class.  A defined constant whose definition may fail is not
  encoded as an atomic constant.
\item[(S3)] The rows of $\mathrm{FUN\text{-}DEF}$ are the eliminations
  of the kernel's defined heads.  The hypotheses H1--H6 of
  Chapter~\ref{ch:meta-functoids} hold.
\item[(S4)] Every predicate of $T$ without a row of
  $\mathrm{PRED\text{-}DEF}$ has a predicative definition, so the
  elimination lemma reads through it.
\item[(S5)] $T$ consists of the catalog and of installed theorems.  The
  argument covers the axioms of the catalog and the definitions read
  as above.  A guarded definition is read as a description
  (Chapter~\ref{ch:meta-functoids}, \DS{31}).
\end{enumerate}
(S2)--(S4) are obligations of the translator, K3
(Section~\ref{sec:meta-open}).  The definitions of the library are
covered by their general form, not one by one.
